\subsection{Công cụ và không gian làm việc}

Nhóm lựa chọn và kết hợp các công cụ hiện đại nhằm tối ưu hóa quá trình giao tiếp, quản lý mã nguồn, soạn thảo tài liệu và theo dõi tiến độ công việc:

\begin{itemize}
    \item \textbf{\LaTeX{} (Công cụ soạn thảo báo cáo học thuật):}
    \begin{itemize}
        \item Sử dụng hệ thống soạn thảo \LaTeX{} để chuẩn hóa cấu trúc báo cáo, tự động hóa việc tạo mục lục, quản lý tài liệu tham khảo và định dạng phông chữ theo chuẩn học thuật.
        \item Cấu trúc báo cáo thành các tệp nguồn (\texttt{.tex}) nhỏ dạng mô-đun, giúp các thành viên dễ dàng phân chia từng phần nội dung để soạn thảo độc lập.
        \item Đảm bảo tệp báo cáo PDF đầu ra đạt chất lượng cao về mặt thẩm mỹ, nhất quán về kiểu chữ và bố cục trang.
    \end{itemize}

    \item \textbf{Slack (Không gian giao tiếp \& Thảo luận):}
    \begin{itemize}
        \item Đóng vai trò là kênh trao đổi trực tiếp hàng ngày giữa các thành viên trong nhóm.
        \item Phân chia các kênh (channels) chuyên biệt như \texttt{\#general} (thông báo chung), \texttt{\#job-description} (thảo luận nội dung bài tập), và \texttt{\#report} (tổng hợp báo cáo) giúp luồng thông tin mạch lạc, tránh trôi tin nhắn.
        \item Tích hợp webhook thông báo tự động từ GitHub để các thành viên kịp thời nắm bắt khi có cập nhật mới.
    \end{itemize}
    
    \item \textbf{GitHub (Quản lý mã nguồn \& Tiến độ):}
    \begin{itemize}
        \item Repository chính thức của nhóm: \url{https://github.com/pnhlongbk2024/SS004.6---Job-Description-26730046---26730026---26730050-}
        \item Lưu trữ và quản lý phiên bản toàn bộ mã nguồn bài tập (các tệp tin cấu trúc \LaTeX{} \texttt{.tex}).
        \item Áp dụng quy trình làm việc chuẩn với Git (Branching, Pull Request, Code Review) nhằm kiểm duyệt nội dung trước khi hợp nhất vào nhánh chính (\texttt{master}).
        \item Hỗ trợ theo dõi lịch sử chỉnh sửa, khôi phục phiên bản khi cần và đồng bộ hóa làm việc nhóm một cách nhất quán.
    \end{itemize}
\end{itemize}
